%-----------------------------------------------------------------------------------------------------------------------------------------------%
%    The MIT License (MIT)
%
%    Copyright (c) 2021 Philip Empl
%
%    Permission is hereby granted, free of charge, to any person obtaining a copy
%    of this software and associated documentation files (the "Software"), to deal
%    in the Software without restriction, including without limitation the rights
%    to use, copy, modify, merge, publish, distribute, sublicense, and/or sell
%    copies of the Software, and to permit persons to whom the Software is
%    furnished to do so, subject to the following conditions:
%    
%    THE SOFTWARE IS PROVIDED "AS IS", WITHOUT WARRANTY OF ANY KIND, EXPRESS OR
%    IMPLIED, INCLUDING BUT NOT LIMITED TO THE WARRANTIES OF MERCHANTABILITY,
%    FITNESS FOR A PARTICULAR PURPOSE AND NONINFRINGEMENT. IN NO EVENT SHALL THE
%    AUTHORS OR COPYRIGHT HOLDERS BE LIABLE FOR ANY CLAIM, DAMAGES OR OTHER
%    LIABILITY, WHETHER IN AN ACTION OF CONTRACT, TORT OR OTHERWISE, ARISING FROM,
%    OUT OF OR IN CONNECTION WITH THE SOFTWARE OR THE USE OR OTHER DEALINGS IN
%    THE SOFTWARE.
%    
%
%-----------------------------------------------------------------------------------------------------------------------------------------------%

%============================================================================%
%
%    DOCUMENT DEFINITION
%
%============================================================================%

\documentclass[10pt,A4,english]{article}    

%----------------------------------------------------------------------------------------
%    ENCODING
%----------------------------------------------------------------------------------------

\usepackage[utf8]{inputenc}        
\usepackage[USenglish]{isodate}
\usepackage{fancyhdr}
\usepackage[numbers]{natbib}

%----------------------------------------------------------------------------------------
%    LOGIC
%----------------------------------------------------------------------------------------

\usepackage{xstring, xifthen}
\usepackage{enumitem}
\usepackage[polish]{babel}
\usepackage{blindtext}
\usepackage{pdfpages}
\usepackage{changepage}

%----------------------------------------------------------------------------------------
%    FONT BASICS
%----------------------------------------------------------------------------------------

\usepackage[default]{raleway}

\renewcommand*\familydefault{\sfdefault}     
\usepackage[T1]{fontenc}
\usepackage{moresize}

%----------------------------------------------------------------------------------------
%    FONT AWESOME ICONS
%---------------------------------------------------------------------------------------- 

\usepackage{fontawesome}

\newcommand{\vcenteredinclude}[1]{\begingroup
\setbox0=\hbox{\includegraphics{#1}}%
\parbox{\wd0}{\box0}\endgroup}
\newcommand{\tab}[1]{\hspace{.2\textwidth}\rlap{#1}}
\newcommand*{\vcenteredhbox}[1]{\begingroup
\setbox0=\hbox{#1}\parbox{\wd0}{\box0}\endgroup}

\newcommand{\icon}[3] {                              
    \makebox(#2, #2){\textcolor{maincol}{\csname fa#1\endcsname}}
}    

\newcommand{\icontext}[4]{                         
    \vcenteredhbox{\icon{#1}{#2}{#3}}  \hspace{2pt}  \parbox{0.9\mpwidth}{\textcolor{#4}{#3}}
}

\newcommand{\iconhref}[5]{                         
    \vcenteredhbox{\icon{#1}{#2}{#5}}  \hspace{2pt} \href{#4}{\textcolor{#5}{#3}}
}

\newcommand{\iconemail}[5]{                         
    \vcenteredhbox{\icon{#1}{#2}{#5}}  \hspace{2pt} \href{mailto:#4}{\textcolor{#5}{#3}}
}

%----------------------------------------------------------------------------------------
%    PAGE LAYOUT  DEFINITIONS
%----------------------------------------------------------------------------------------

\usepackage{paracol}
\usepackage{tikzpagenodes}
\usetikzlibrary{calc}
\usepackage{lmodern}
\usepackage{multicol}
\usepackage{lipsum}
\usepackage{atbegshi}
\usepackage[a4paper]{geometry}

\geometry{top=1cm, bottom=1cm, left=1cm, right=1cm}

\usepackage{fancyhdr}
\pagestyle{empty}

\setlength{\parindent}{0mm}

%----------------------------------------------------------------------------------------
%    TABLE /ARRAY DEFINITIONS
%---------------------------------------------------------------------------------------- 

\usepackage{array}

\newcolumntype{x}[1]{%
>{\raggedleft\hspace{0pt}}p{#1}}%

%----------------------------------------------------------------------------------------
%    GRAPHICS DEFINITIONS
%---------------------------------------------------------------------------------------- 

\usepackage{graphicx}
\usepackage{tikz}            
\usepackage{ragged2e}    
\usetikzlibrary{shapes, backgrounds,mindmap, trees}

%----------------------------------------------------------------------------------------
%    Color DEFINITIONS
%---------------------------------------------------------------------------------------- 
\usepackage{transparent}
\usepackage{color}

\definecolor{maincol}{RGB}{ 64,64,64}
\definecolor{darkcol}{RGB}{ 70, 70, 70 }
\definecolor{lightcol}{RGB}{245,245,245}
\definecolor{accentcol}{RGB}{59,77,97}

\usepackage[hidelinks]{hyperref}

\newcommand{\mpwidth}{\linewidth-\fboxsep-\fboxsep}
    
%============================================================================%
%
%    CV COMMANDS
%
%============================================================================%

\newcommand{\cvlist}[1] {
    \begin{itemize}{#1}\end{itemize}
}

\newcommand{\cvtext}[1] {
    \begin{tabular*}{1\mpwidth}{p{0.98\mpwidth}}
        \parbox{1\mpwidth}{#1}
    \end{tabular*}
}
\newcommand{\cvtextsmall}[1] {
    \begin{tabular*}{0.8\mpwidth}{p{0.8\mpwidth}}
        \parbox{0.8\mpwidth}{#1}
    \end{tabular*}
}

\newcommand{\cvsection}[1] {
    \vspace{14pt}
    \cvtext{
        \textbf{\LARGE{\textcolor{darkcol}{#1}}}\\[-4pt]
        \textcolor{accentcol}{ \rule{0.2\textwidth}{1.5pt} } \\
    }
}

\newcommand{\cvsectionsmall}[1] {
    \vspace{14pt}
    \cvtext{
        \textbf{\Large{\textcolor{darkcol}{#1}}}\\[-4pt]
        \textcolor{accentcol}{ \rule{0.2\textwidth}{1.5pt} } \\
    }
}

\newcommand{\cvheadline}[1] {
    \vspace{16pt}
    \cvtext{
        \textbf{\Huge{\textcolor{accentcol}{#1}}}\\[-4pt]
         
    }
}

\newcommand{\cvsubheadline}[1] {
    \vspace{16pt}
    \cvtext{
        \textbf{\huge{\textcolor{darkcol}{#1}}}\\[-4pt]
         
    }
}

\newcommand{\cvskill}[3] {
    \begin{tabular*}{1\mpwidth}{p{0.72\mpwidth}  r}
         \textcolor{black}{\textbf{#1}} & \textcolor{maincol}{#2}\\
    \end{tabular*}%
    
    \hspace{4pt}
    \begin{tikzpicture}[scale=1,rounded corners=2pt,very thin]
        \fill [lightcol] (0,0) rectangle (1\mpwidth, 0.15);
        \fill [accentcol] (0,0) rectangle (#3\mpwidth, 0.15);
      \end{tikzpicture}%
}

\newcommand{\cvevent}[7] {
    \parbox{\mpwidth}{
        \begin{tabular*}{1\mpwidth}{p{0.66\mpwidth}  r}
             \textcolor{black}{\textbf{#2}} & \colorbox{accentcol}{\makebox[0.32\mpwidth]{\textcolor{white}{\textbf{#1}}}} \\
            \textcolor{maincol}{#3} & \\
        \end{tabular*}\\[8pt]
    
        \ifthenelse{\isempty{#4}}{}{
            \cvtext{\raggedright #4}\\
        }
    }
    \vspace{14pt}
}

\newcommand{\cvmetaevent}[4] {
    \textcolor{maincol} { \cvtext{\textbf{\begin{flushleft}#1\end{flushleft}}}}

    \ifthenelse{\isempty{#2}}{}{
    \textcolor{black} {\cvtext{\textbf{#2}} }
    }

    \ifthenelse{\isempty{#3}}{}{
        \cvtext{{ \textcolor{maincol} {#3} }}\\
    }

    \cvtext{\raggedright #4}\\[14pt]
}

\newcommand{\cvqrcode}[1] {
    \begin{center}
        \includegraphics[width={#1}\mpwidth]{qrcode}
    \end{center}
}

% HEADER AND FOOOTER 
%====================================
\newcommand\Header[1]{%
\begin{tikzpicture}[remember picture,overlay]
\fill[accentcol]
  (current page.north west) -- (current page.north east) --
  ([yshift=50pt]current page.north east|-current page text area.north east) --
  ([yshift=50pt,xshift=-3cm]current page.north|-current page text area.north) --
  ([yshift=10pt,xshift=-5cm]current page.north|-current page text area.north) --
  ([yshift=10pt]current page.north west|-current page text area.north west) -- cycle;
\node[font=\sffamily\bfseries\color{white},anchor=west,
  xshift=0.7cm,yshift=-0.32cm] at (current page.north west)
  {\fontsize{12}{12}\selectfont {#1}};
\end{tikzpicture}%
}

\newcommand\Footer[1]{%
\begin{tikzpicture}[remember picture,overlay]
\fill[lightcol]
  (current page.south east) -- (current page.south west) --
  ([yshift=-80pt]current page.south east|-current page text area.south east) --
  ([yshift=-80pt,xshift=-6cm]current page.south|-current page text area.south) --
  ([xshift=-2.5cm,yshift=-10pt]current page.south|-current page text area.south) --    
  ([yshift=-10pt]current page.south east|-current page text area.south east) -- cycle;
\node[yshift=0.32cm,xshift=9cm] at (current page.south) {\fontsize{10}{10}\selectfont \textbf{\thepage}};
\end{tikzpicture}%
}

%=====================================
%    DOCUMENT CONTENT
%============================================================================%
\begin{document}

\columnratio{0.31}
\setlength{\columnsep}{2.2em}
\setlength{\columnseprule}{4pt}
\colseprulecolor{white}


\AtBeginShipoutFirst{\Header{CV}\Footer{1}}
\AtBeginShipout{\AtBeginShipoutAddToBox{\Header{CV}\Footer{2}}}

\newpage

\colseprulecolor{lightcol}
\columnratio{0.31}
\setlength{\columnsep}{2.2em}
\setlength{\columnseprule}{4pt}
\begin{paracol}{2}


\begin{leftcolumn}
%---------------------------------------------------------------------------------------
%    META IMAGE
%----------------------------------------------------------------------------------------
\includegraphics[width=\linewidth]{resources/house.jpg}    

%---------------------------------------------------------------------------------------
%    META SKILLS
%----------------------------------------------------------------------------------------
    \fcolorbox{white}{white}{\begin{minipage}[c][1.5cm][c]{1\mpwidth}
        \LARGE{\textbf{\textcolor{maincol}{Gregory House}}} \\[2pt]
        \normalsize{ \textcolor{maincol} {Lekarz Medycyny (M.D.)} }
\end{minipage}} \\
\icontext{CaretRight}{12}{11.06.1959 r. w USA}{black}\\[6pt]
\icontext{CaretRight}{12}{Obywatel USA}{black}\\[6pt]
\icontext{CaretRight}{12}{Kawaler}{black}\\[6pt]



\cvsection{Umiejętności}

\cvskill{Diagnostyka różnicowa} {25+ lat} {1} \\[-2pt]

\cvskill{Choroby zakaźne} {20+ lat} {0.95} \\[-2pt]

\cvskill{Nefrologia} {20+ lat} {0.95} \\[-2pt]

\cvskill{Farmakologia i toksyk.} {20+ lat} {0.90} \\[-2pt]

\cvskill{Logika i dedukcja} {Ekspert} {1} \\[-2pt]

\cvskill{Gra na instr.} {Hobby} {0.80} \\[-2pt] \\

\cvskill{Angielski} {Native} {1} \\[-2pt]

\cvskill{Polski} {C1} {0.85} \\[-2pt]

\cvskill{Mandaryński / Hindi} {Podstawy} {0.60} \\[-2pt]

%---------------------------------------------------------------------------------------
%    EDUCATION
%----------------------------------------------------------------------------------------
\newpage
\cvsection{Edukacja}

\cvmetaevent
{1984}
{Doktor Nauk Med. (M.D.)}
{University of Michigan Medical School}
{\textit{Dyplom z wyróżnieniem.} \newline Zmiana uczelni po incydencie na kampusie Hopkinsa.}

\cvmetaevent
{1979 - 1982}
{Studia medyczne}
{Johns Hopkins University}
{\textit{Medycyna} \newline Wydalony za oszukiwanie, którego nie żałuję, bo miałem rację.}

\cvmetaevent
{1979}
{B.S. w dziedzinie fizyki}
{Johns Hopkins University}
{\textit{Studia przedmedyczne} \newline Zrozumienie fundamentalnych praw natury.}


\cvsection{Zainteresowania}

\icontext{CaretRight}{12}{Gra na instrumentach}{black}\\[6pt]
\icontext{CaretRight}{12}{Motocykle (Honda CBR1000RR)}{black}\\[6pt]
\icontext{CaretRight}{12}{Opery mydlane}{black}\\[6pt]
\icontext{CaretRight}{12}{Medycyna niekonwencjonalna}{black}\\[6pt]
\icontext{CaretRight}{12}{Udowadnianie innym błędów}{black}\\[6pt]

\end{leftcolumn}
\begin{leftcolumn}

\cvsection{Kontakt}

\icontext{MapMarker}{16}{Wydział Lekarski\\ul. Mikulicza-Radeckiego 5\\50-345 Wrocław}{black}\\[6pt]
\icontext{MobilePhone}{16}{+48 987 654 321}{black}\\[6pt]
\iconemail{Envelope}{16}{g.house@umed.wroc.pl}{g.house@umed.wroc.pl}{black}\\[6pt]
\iconhref{Home}{16}{princeton-plainsboro.edu}{http://www.princeton-plainsboro.edu}{black}\\[6pt]

\end{leftcolumn}
\begin{rightcolumn}
%---------------------------------------------------------------------------------------
%    PROFILE
%----------------------------------------------------------------------------------------
\cvsection{Profil Zawodowy}
\vspace{4pt}

\cvtext{\raggedright
Jestem wybitnym diagnostą z podwójną specjalizacją: z nefrologii oraz chorób zakaźnych. Przez ponad dekadę z sukcesami kierowałem Oddziałem Medycyny Diagnostycznej w Princeton-Plainsboro Teaching Hospital, rozwiązując najtrudniejsze przypadki medyczne, z którymi reszta amerykańskiego środowiska lekarskiego nie potrafiła sobie poradzić. 

Moja filozofia dydaktyczna jest prosta: ,,Wszyscy kłamią, ale objawy nie''. Aplikuję na stanowisko wykładowcy na Uniwersytecie Medycznym we Wrocławiu, ponieważ dostrzegam palącą potrzebę wykształcenia nowej generacji lekarzy. Lekarzy, którzy potrafią myśleć logicznie i poza schematami, zamiast bezmyślnie odhaczać procedury z podręcznika i ślepo ufać wynikom z laboratorium. Oferuję uczelni mój analityczny rygor, bezwzględną szczerość i zmysł obserwacji, którego nie da się wyuczyć z książek.
}


%---------------------------------------------------------------------------------------
%    WORK EXPERIENCE
%----------------------------------------------------------------------------------------

\vspace{10pt}
\cvsection{Doświadczenie zawodowe}
\vspace{4pt}


\cvevent
{2004 - 2012}
    {Ordynator Oddziału Medycyny Diagnostycznej}
    {Princeton-Plainsboro Teaching Hospital \newline Princeton, New Jersey, USA}
    {Kierowanie elitarnym zespołem diagnostycznym (rezydentów i specjalistów). Selekcja i rozwiązywanie wyłącznie najrzadszych, najbardziej skomplikowanych przypadków medycznych. Stosowanie nieszablonowych metod diagnostycznych (w tym nieautoryzowanych przeszukań domów pacjentów w poszukiwaniu toksyn). Prowadzenie cotygodniowych wykładów klinicznych dla studentów medycyny i stażystów.}
    \vfill\null


\cvevent
{2012 - dzisiaj}
    {Niezależny Konsultant Medyczny i Diagnosta}
    {Praktyka prywatna / Uczelnie medyczne \newline Cały świat (w tym Europa i Azja)}
    {Udzielanie konsultacji w beznadziejnych przypadkach klinicznych dla szpitali akademickich. Prowadzenie gościnnych wykładów z zakresu dedukcji medycznej i zaawansowanej diagnostyki różnicowej. Badania nad niestandardowymi terapiami zakażeń opornych na antybiotyki.}
    \vfill\null


    
\cvevent
{1990 - 2004}
    {Specjalista ds. Chorób Zakaźnych i Nefrologii}
    {Różne ośrodki kliniczne \newline USA}
    {Uzyskanie podwójnej specjalizacji, prowadzenie pionierskich badań nad nietypowymi objawami chorób autoimmunologicznych (udowadnianie rezydentom, że to rzadko jest toczeń). Wykłady z anatomii patologicznej i farmakologii klinicznej.}
    \vfill\null


\end{rightcolumn}

\colseprulecolor{lightcol}
\columnratio{0.31}
\setlength{\columnsep}{2.2em}
\setlength{\columnseprule}{4pt}

\begin{rightcolumn}
\newpage
\cvsection{Wybrane publikacje}

\begin{itemize}[leftmargin=*]
\item House, G. (2010). "Prawda w diagnostyce: Dlaczego wywiad medyczny to fikcja literacka". W: \textit{Journal of Medical Diagnostics}, 12(4), 112-145.
\item House, G. \& Wilson, J. (2008). "Onkologia a choroby zakaźne: Granice błędów poznawczych w diagnozach różnicowych". W: \textit{New England Journal of Medicine}, 358(10).
\item House, G. (2006). "Lupus: Kiedy to nie jest toczeń (czyli prawie zawsze)". W: \textit{Annals of Internal Medicine}, 145(2).
\item House, G. (1998). "Mechanizmy oporności na leki w rzadkich infekcjach grzybiczych - analiza przypadków beznadziejnych". W: \textit{Clinical Infectious Diseases}, 26(3).
\end{itemize}

% hofixes to create fake-space to ensure the whole height is used
\mbox{}
\vfill
\mbox{}
\vfill
\mbox{}
\vfill
\mbox{}
\vfill
\mbox{}
\vfill
\mbox{}
\vfill
\mbox{}
\vfill
\mbox{}

Wrocław, \today     \hspace{1cm}   \hrulefill

\hspace*{30mm}\phantom{Wrocław, \today }dr n. med. Gregory House

\end{rightcolumn}
\end{paracol}


\end{document}